\section{Описание конечного состояния}
\label{sec:endpoint}

\subsection{Восстановление по суффиксу}

Пусть состояние ученика принадлежит инвариантному множеству
$\mathcal S\subset\mathbb R^d$ и обновляется по метке из конечного алфавита
$\mathcal Y$:
$\theta_{t+1}=F_{Y_t}(\theta_t)$.
Декодер стартует из известного $\theta_0\in\mathcal S$. Требуется восстановить
предиктор $P_{\theta_T}$ с заданной ошибкой, а не обязательно повторить
все $T$ обновлений. Разрешены вспомогательные учителя, детерминированно
выбирающие нужную метку.

\begin{theorem}[Длина описания сжимающего ученика]
\label{thm:suffix}
Пусть для всех $y\in\mathcal Y$ и $\theta,\theta'\in\mathcal S$
\begin{align*}
 \|F_y(\theta)-F_y(\theta')\|&\le\rho\|\theta-\theta'\|,
                  &&0<\rho<1,\\
 \KL(P_\theta\Vert P_{\theta'})&\le C\|\theta-\theta'\|^2,
                  &&0<C<\infty,\\
 P_\theta(y)&\ge p_{\min}>0.
\end{align*}
Предположим также, что $\sup_{\theta\in\mathcal S}\|\theta-\theta_0\|\le R$
при некотором $R>0$.
Для $\delta>0$ положим
\begin{equation}
 h_\delta=\max\left\{0,
 \left\lceil\frac{\ln(CR^2/\delta)}{2\ln(1/\rho)}\right\rceil\right\},
 \qquad h=\min\{T,h_\delta\}.
 \label{eq:suffix-length}
\end{equation}
Повторение последних $h$ меток исходной траектории из $\theta_0$ даёт
состояние $\theta'$ с
$\KL(P_{\theta_T}\Vert P_{\theta'})\le\delta$.
Если для сокращённой записи используется заранее согласованный поток
кандидатов, независимый от исходной траектории, то существует сообщение
с ожидаемой длиной не выше
\begin{equation}
 \ell_\delta(h+1)+hL(-\ln p_{\min}).
 \label{eq:endpoint-code}
\end{equation}
В частности, равномерно по $T$ ожидаемая длина ограничена величиной
\begin{equation}
 \ell_\delta(h_\delta+1)+h_\delta L(-\ln p_{\min})
   =O(1+\ln(1/\delta)),\qquad \delta\downarrow0,
 \label{eq:endpoint-uniform}
\end{equation}
при фиксированных $C,R,\rho,p_{\min}$.
\end{theorem}
\begin{proof}
Если $T<h_\delta$, повторяем всю последовательность и получаем точное
состояние. Иначе истинное состояние перед суффиксом ---
$\theta_{T-h}$ --- удалено от $\theta_0$ не более чем на $R$.
Применение одной и той же последовательности из $h$ сжимающих отображений
даёт $\|\theta_T-\theta'\|\le R\rho^h$.
Тогда дивергенция не выше $CR^2\rho^{2h}\le\delta$.
Случай $h_\delta=0$ соответствует $CR^2\le\delta$.

Заголовок --- $\delta$-код числа $h+1$ --- сообщает декодеру количество
меток. Каждую нужную метку $y$ выбирает вспомогательный учитель с точечной
мерой $\delta_y$. Передаётся индекс её первого появления среди
независимых кандидатов из текущего предиктора. Для фиксированной исходной
траектории такой индекс имеет распределение
$\Geom(P_{\theta'}(y))$, и его ожидаемая стоимость не выше
$L(-\ln p_{\min})$. Независимость нового потока от исходной траектории
обеспечивает эту условную оценку даже при выборе суффикса по исходным данным.
После передачи каждой метки состояние воспроизводится точно. Складывая
стоимости, получаем~\eqref{eq:endpoint-code}.
Наконец, $h\le h_\delta$, функция $\ell_\delta$ не убывает,
$h_\delta=O(\ln(1/\delta))$, а длина заголовка имеет порядок
$O(1+\ln(h_\delta+1))$. Это доказывает~\eqref{eq:endpoint-uniform}.
\end{proof}

\begin{corollary}[Линейный разрыв с полной записью]
\label{cor:gap}
При условиях теоремы~\ref{thm:suffix} и фиксированной $\delta>0$ полная
requential-запись из $T$ единичных пакетов с отдельным индексом на шаг
имеет длину не менее $T$ битов. Минимальная ожидаемая стоимость описания
конечной точки ограничена константой, не зависящей от $T$.
\end{corollary}
\begin{proof}
Каждый индекс стоит хотя бы один бит. Верхняя оценка для конечной точки
дана~\eqref{eq:endpoint-code}.
\end{proof}

Это сравнение меняет цель декодирования: короткое сообщение восстанавливает
конечный предиктор с допуском, но не исходную случайную траекторию. Поэтому
разрыв не противоречит оптимальности исходного кода для более строгой задачи.
Предположение об инвариантности и сжатии существенно: для накопительного
среднего старые метки не теряют относительный вес экспоненциально, и
фиксированный суффикс обычно недостаточен.

\begin{corollary}[Аффинный бернуллиевский ученик]
\label{cor:affine}
Для обновления
\begin{equation}
 F_y(p)=\frac{p+1/8+3y/4}{2},\qquad
 p_0=\frac12,\qquad y\in\{0,1\},
 \label{eq:affine}
\end{equation}
отрезок $[1/8,7/8]$ инвариантен. Если $T\ge h$ и $p'$ восстановлено из
$p_0$ по последним $h$ меткам, то
\begin{equation}
 |p_T-p'|\le\frac38\,2^{-h},\qquad
 \KL(\Bern(p_T)\Vert\Bern(p'))\le\frac97\,4^{-h}.
 \label{eq:affine-bound}
\end{equation}
При $h=6$ дивергенция меньше $10^{-3}$ ната для любой исходной траектории.
\end{corollary}
\begin{proof}
Непосредственная проверка концов отрезка доказывает инвариантность.
Каждое обновление уменьшает разность состояний вдвое; начальное расхождение
между $p_{T-h}$ и $1/2$ не выше $3/8$.
Из $\ln x\le x-1$ следует
\[
 \KL(\Bern(p)\Vert\Bern(p'))
 \le\frac{(p-p')^2}{p'(1-p')}.
\]
Знаменатель не меньше $7/64$. Подстановка первой оценки
даёт~\eqref{eq:affine-bound}, причём $(9/7)4^{-6}<10^{-3}$.
\end{proof}

\subsection{Когда кодовая длина даёт гарантию обобщения}

Связь длины описания предиктора с риском на новых данных задаётся стандартной
границей для префиксного кода \citep{grunwald2004,arora2018}.

\begin{proposition}[Стандартная граница для кода предиктора]
\label{prop:occam}
Пусть $Z_1,\dots,Z_n$ независимы и одинаково распределены, убыток лежит в
$[0,1]$, а префиксный код и декодер $m\mapsto h_m$ фиксированы независимо от
выборки. Обозначим популяционный и эмпирический риски через
$R(h)=\E\mathcal L(h,Z)$ и
$\widehat R_n(h)=n^{-1}\sum_{i=1}^n\mathcal L(h,Z_i)$.
Для любого $\eta\in(0,1)$ с вероятностью не меньше $1-\eta$ одновременно для всех
сообщений $m$
\begin{equation}
 R(h_m)\le\widehat R_n(h_m)
 +\sqrt{\frac{\ell(m)\ln2+\ln(1/\eta)}{2n}}.
 \label{eq:occam}
\end{equation}
Если декодер использует общую случайность, она также должна быть независима
от выборки. В границу подставляется реализованная длина полного сообщения.
\end{proposition}
\begin{proof}
Условимся на общей случайности. Для каждого фиксированного сообщения
неравенство Хёффдинга ограничивает вероятность нарушения
\eqref{eq:occam} величиной $\eta2^{-\ell(m)}$.
По неравенству Крафта и объединению событий сумма этих вероятностей не
превосходит $\eta$. Остаётся усреднить по общей случайности.
\end{proof}

Сообщение $m$ может выбираться по данным, поскольку граница одновременна.
Длина $\ell(m)$ включает заголовок и все выбранные по данным параметры
декодирования. Граница относится к ограниченной функции потерь и к
декодированному предиктору $h_m$.
